\begin{figure}
    \centering
    \includegraphics[width=0.99\textwidth]{content/figures/hdcells.pdf}
    \caption{\textbf{Individual differences in head-direction coding around a shared ring manifold}. \textbf{a.} Left, example neurons. Right, all neurons sorted (cross-validated across time) by their preferred direction angle. \textbf{b.} Left, aligning all subjects (n=31) to a common reference (Methods) reveals a high-dimensional ring. Bottom, the first principal components of the aligned space (dark blue) explain a substantial amount of variance of the original space (light blue), suggesting individual differences. Right, illustration of a zero-shot learning decoding procedure from 2 aligned animals. As we aligned a larger number of subjects, the decoder approached the performance of the decoder trained and tested on independent folds of the same animal. \textbf{c.} Left, categoricality of head-direction tuning curves within each subject (as in \textbf{Fig.}~2b): \hl{XX} of 31 subjects exceed $|z| = 1.96$, i.e.\ tuning curves are not randomly mixed within animals. Middle, neurons pooled across subjects are also categorical: the mean silhouette (red) exceeds a null in which all neurons form one continuous cloud. Right, subjects themselves do not cluster in shape space: the best silhouette across subjects (red) does not exceed the one-blob null (blue; $p = 0.23$). \textbf{d.} Left, illustration of the cross-validation approach. Right, independent folds of animal A (blue and purple, top) have smaller mean squared-error (MSE) distance after alignment than a similar analysis between animal A and B (blue and red, below). \textbf{e.} Same analyses as in d), repeated for all animals and 1 individual split. Distances between two different splits of the same animal are significantly smaller than across different animals (t=-3.13, p=0.0017). Middle, the former replicated on n=100 different splits. Whitening the data before alignment captures individual differences slightly better (p=0.05). MDS of the dissimilarity matrix suggests that tuning width drives across-animal differences (red gradient vs green gradient). Inset, decoding $R^2$ (using bi-cross validation in b) of tuning width from this space is significantly above chance (95\% C.I. of shuffle decoder in gray). \textbf{f.} Analyses on a separate dataset of n=6 animals shows that a rotation in the cue space (90\textdegree) is reflected in a consistent angle rotation in the neural space (98\textdegree).}
    \label{fig:hdcells}
\end{figure}

\subsection{Animal-to-animal variability in head direction coding}

Having found that population geometry is organized across brain regions, we now turn to another scale: variability across individuals. Specifically, we ask whether the geometry of a single region differs systematically from one animal to the next. With this aim, we use shape distances to illuminate animal-to-animal variability in two recently published datasets of head direction cells (\cite{duszkiewicz}).
In both datasets, spiking activity in the mouse postsubiculum (PoSub) was collected using silicon probes during open-field exploration while position and heading direction were captured by an overhead video camera (for details, see \cite{duszkiewicz}).
The first dataset consisted of $K = 31$ sessions / subjects and a total of $2691$ neurons.
We first collected single-unit tuning curves for heading across all sessions and sorted them according to peak activity.
Preferred heading directions were roughly uniformly distributed (\textbf{Fig.} \ref{fig:hdcells}a) and applying \textit{PCA} to this aggregated dataset revealed a high-dimensional ring structure (\hl{sup}), consistent with the attractor manifold previously described in this circuit \parencite{chaudhuri2019intrinsic} and akin to the simulated dataset we analyzed above. 

To confirm that animals shared a common manifold, we attemp to decode head direction in one heldout subject from training data recorded in other subjects.
We started by using rotation and reflections to align all subjects to a common reference---a procedure known as \textit{generalized Procrustes analysis} \parencite{gower1975generalized}.
This revealed a high-dimensional ring structure (\textbf{Fig.} \ref{fig:hdcells}b), similar to performing PCA on all data (\hl{sup}) confirming that this structure is shared by all subjects and not just present in a subset of them. Because all animals are aligned to a common shape (in this case, a ring), we could train a decoder on a subset of animals and predict head direction angles on left-out subjects, essentially performing zero-shot learning (\textbf{Fig.} \ref{fig:hdcells}b, right). An analogous strategy has been used to decode movement across 4 individuals from aligned latent dynamics in motor cortex \parencite{safaie2023preserved}. 

We noted however that the first principal components of the aligned space explained significantly less variance than the original space of each animal (\textbf{Fig.} \ref{fig:hdcells}b, bottom left), suggesting individual differences despite a common manifold across subjects. To our knowledge, it is unknown whether there are systematic animal-to-animal differences in the statistics of head direction tuning curves in PoSub. To investigate this, we first asked whether single neuron selectivity were randomly mixed within each subject, applying the clustering analysis that \textcite{posani2026} developed to characterize single-neuron selectivity across cortex. Unsurprisingly, head direction tuning curves in PoSub were not randomly mixed (\textbf{Fig.} \ref{fig:hdcells}c), unlike the largely continuous selectivity profiles that method reports in cortical areas \textcite{posani2026}. We next asked our main question: are individual differences random or organized in subject space? We answered this question by representing each animal by its Procrustes distances to every other and found that subjects did not clustered more than expected under a null (\textbf{Fig.} \ref{fig:hdcells}c, right). As when characterizing strcuture across regions (\textbf{Fig. }\ref{fig:allen}), this does not imply that idiosyncrasies in each subject's geometry is unstructured. To clarify this, we thus measured the Procrustes distance between two disjoint partitions of the data along neuron and time dimensions (see illustration of bi-cross validation in \textbf{Fig.} \ref{fig:hdcells}d, left; \textit{Methods}). On average, the distance between partitions of the same animal was smaller than the distance between partitions taken from different animals (\textbf{Fig.} \ref{fig:hdcells}e). To investigate what features of the neural tuning curves account for these individual differences in ring shape, we first used unsupervised methods to perform an exploratory analysis.
As before, we used MDS followed by PCA to visualize each subject within a 2D space. Coloring each session by the average width and average amplitude of the head direction tuning curves suggested that the former, and not the latter, was a principal driver of inter-animal differences (\textbf{Fig.} \ref{fig:hdcells}e, right). 
To confirm this finding in our exploratory analysis, we turned to supervised methods.
The decoding of the tuning width from this space was significantly above chance (bi-cross validated as in \textbf{Fig.} \ref{fig:hdcells}d), demonstrating that indeed tuning width is driving the inter-individual differences captured by the shape metrics (\textbf{Fig.} \ref{fig:hdcells}e, inset). We found similar signatures of individual differences in a second, smaller dataset ($K = 6$ animals).
In this experiment, open exploration was performed in boxes with a salient landmark which was rotated by 90 degrees during the course of the recording (cue 1 vs cue 2, \textbf{Fig.} \ref{fig:hdcells}f).
We find that the optimal rotation identified by Procrustes analysis in neural firing rate space was roughly matched in magnitude to the cue rotation (\textbf{Fig.} \ref{fig:hdcells}f).

In sum, all subjects encode the head direction angle in a ring, but each of these rings has idiosyncratic distortions in higher dimensions. Furthermore, we found that this distortions are organized by tuning curve proprieties by performing regression in subject space. 

  

